%% ---------------------------------------------------------------------------
\section{The cost of a post-quantum content credential}
\label{sec:cost}

\begin{table}[t]
\caption{Manifest cost by signature algorithm with the signer's certificate and its intermediate in \texttt{x5chain}, the root omitted as C2PA~2.5 says it should be. Sizes are exact; the percentage is the median over the Kodak set with its range. \emph{vs ES256} and \emph{vs PS256} compare with the two algorithms deployed signers use (Section~\ref{sec:prodchains}). Signing time is the signature over the COSE \texttt{Sig\_structure} on Apple M1 Pro, through liboqs for ML-DSA and OpenSSL for the classical rows. Times are medians over 10 independent runs, each the median of its own repetitions; Table~\ref{tab:ci} in the supplement gives their intervals across runs.}
\label{tab:pq}
\small
\setlength{\tabcolsep}{4pt}
\begin{tabular}{lrrrrrrrr}
\toprule
 & COSE & Sig. & Chain & Manifest & \% of image & vs & vs & Sign\\
Algorithm & id & (B) & (B) & (B) & median (range) & ES256 & PS256 & (ms)\\
\midrule
Ed25519 & -8 & 64 & 1,387 & 2,544 & 2.37 (1.35--3.29) & 1.0$\times$ & 0.7$\times$ & 0.030\\
ES256 & -7 & 64 & 1,501 & 2,658 & 2.48 (1.41--3.44) & 1.0$\times$ & 0.7$\times$ & 0.022\\
PS256 & -37 & 256 & 2,291 & 3,642 & 3.39 (1.94--4.71) & 1.4$\times$ & 1.0$\times$ & 0.535\\
ML-DSA-44 & -48 & 2,420 & 8,723 & 12,238 & 11.40 (6.51--15.83) & 4.6$\times$ & 3.4$\times$ & 0.074\\
ML-DSA-65 & -49 & 3,309 & 11,781 & 16,185 & 15.08 (8.61--20.93) & 6.1$\times$ & 4.4$\times$ & 0.116\\
ML-DSA-87 & -50 & 4,627 & 15,697 & 21,419 & 19.95 (11.39--27.70) & 8.1$\times$ & 5.9$\times$ & 0.165\\
\bottomrule
\end{tabular}
\end{table}


Table~\ref{tab:pq} is the central result, and it has to be read against what
signers use today, which is ES256 and, about half as often, RSA-2048 PS256
(Section~\ref{sec:prodchains}). Against ES256 an ML-DSA manifest is
\NvsESDSAfortyfour{}, \NvsESDSAsixtyfive{} and \NvsESDSAeightyseven{} times
larger; against PS256, \NvsPSDSAfortyfour{}, \NvsPSDSAsixtyfive{} and
\NvsPSDSAeightyseven{} times. Either way the manifest grows far more slowly than
the signature, which grows from 64 bytes under ES256 and 256 under RSA-2048 to
3,309 under ML-DSA-65.

The chain column says why. Leaf and intermediate together are
\NchainBytesDSAsixtyfive{} bytes under ML-DSA-65, \NchainShareDSAsixtyfive\% of the signed COSE structure and \NchainShareManDSAsixtyfive\% of the whole manifest once JUMBF framing and the claim are added, against a 3,309-byte signature; under ES256 the
same two certificates are \NchainShareES\% of a much smaller structure. Each
certificate carries a public key and its issuer's signature, so every
certificate in the chain repeats the growth that the manifest signature
shows once. Post-quantum cryptography did not create this imbalance; it scaled
an existing one until it became the cost.

Signing time is not part of the argument. Through liboqs, ML-DSA-44 signs in
\NsignMsDSAfortyfour\,ms, against \NsignMsEd\,ms for Ed25519 and
\NsignMsPS\,ms for PS256 through OpenSSL on the same machine.

\subsection{The denominator problem}
\label{sec:denominator}

A percentage of a Kodak image invites a division the reader will do anyway.
Kodak images are 768$\times$512 and a median of about 108\,kB; a phone
photograph is closer to 4\,MB, where the same ML-DSA-65 manifest that is
\NpctimgDSAsixtyfive\% of a Kodak image is under half a per cent. Once
thumbnails are disabled, manifest size does not depend on the pixels: it is
fixed by the algorithm, by how much of the chain travels, and by the
assertions the generator writes, which for a capture are a handful of hashes
and actions. Ingredients, edit histories and large assertions add bytes of
their own that no choice of algorithm changes. The percentage is therefore an
exact curve, and Figure~\ref{fig:sizeladder} plots it.

\begin{figure}[t]
\centering
\includegraphics[width=0.86\linewidth]{figures/sizeladder}
\Description{Log-log plot of manifest size as a percentage of file size
against file size from 10 kB to 20 MB. Falling straight lines for ML-DSA-65
with timestamp and OCSP, ML-DSA-65 timestamped, ML-DSA-65 with leaf and
intermediate, ML-DSA-65 leaf only, ML-DSA-65 thumbprint, and PS256 timestamped.
A shaded band marks the Kodak images; vertical guides mark a thumbnail, a
social re-encode, a 1080p JPEG and a phone original.}
\caption{Manifest cost against file size. Sizes are measured
(Tables~\ref{tab:timestamp} and~\ref{tab:carriage}) and do not depend on
image content once thumbnails are disabled, so each line is exact. The shaded
band is the Kodak set; the vertical guides are reference sizes, not
measurements. Both axes are logarithmic, so the vertical distance between two
lines is the ratio of their manifest sizes at every file size.}
\label{fig:sizeladder}
\end{figure}

A reader with phone originals in mind will find the migration affordable, and
for that case they are right. A platform storing derivatives, a camera signing
burst frames or a pipeline signing video operates at the left of the figure,
where the manifest is a double-digit percentage of the payload. The two
readings disagree about file size, not about cryptography.

\subsection{The manifest a deployment ships}
\label{sec:timestamp}

\begin{table}[t]
\caption{The manifest a deployment ships: signature, then RFC~3161 timestamp, then stapled revocation evidence, as most production signatures carry a timestamp and a quarter also carry revocation evidence (Section~\ref{sec:prodchains}). Signer and time-stamping authority each carry their certificate and intermediate, roots omitted; the \emph{+ OCSP} column adds a response from a delegated responder, with the responder's certificate inside it, the form the reference implementation honours. Every object is real and signed under the named algorithm. Figure~\ref{fig:compose} breaks the bytes down by role.}
\label{tab:timestamp}
\small
\setlength{\tabcolsep}{4pt}
\begin{tabular}{lrrrrrr}
\toprule
 & Signature & + timestamp & + OCSP & \multicolumn{2}{c}{Timestamped vs} & \% of\\
\cmidrule(lr){5-6}
Algorithm & (B) & (B) & (B) & ES256 & PS256 & image\\
\midrule
Ed25519 & 2,544 & 4,519 & 5,549 & 1.0$\times$ & 0.7$\times$ & 4.2\\
ES256 & 2,657 & 4,722 & 5,819 & 1.0$\times$ & 0.7$\times$ & 4.4\\
PS256 & 3,642 & 6,687 & 8,371 & 1.4$\times$ & 1.0$\times$ & 6.2\\
ML-DSA-44 & 12,238 & 23,915 & 30,977 & 5.1$\times$ & 3.6$\times$ & 22.2\\
ML-DSA-65 & 16,185 & 31,809 & 41,289 & 6.7$\times$ & 4.8$\times$ & 29.5\\
ML-DSA-87 & 21,419 & 42,277 & 55,033 & 9.0$\times$ & 6.3$\times$ & 39.3\\
\bottomrule
\end{tabular}
\end{table}


Deployed manifests carry more than one signature. \NcorpProdTs{} of the
\NcorpProdSigs{} production signatures we read carry an RFC~3161
timestamp~\cite{rfc3161}, which is a second signature with certificates of
its own, and the specification asks for one whenever revocation status is
attached~\cite{c2pacc24}. Revocation evidence is a third signature, a stapled
OCSP response~\cite{rfc6960}, and \NcorpProdOcsp{} of the production
signatures carry it; validators need not check it, the gap that Golaszewski et
al.\ identify~\cite{golaszewski2026}, which makes the last column of
Table~\ref{tab:timestamp} the configuration that closing the gap would make
universal. RFC~6960 lets the issuing CA sign the response itself, and then
it carries no certificate, since the verifier already has the CA's from
\texttt{x5chain}. C2PA~2.5 requires the certificate that signs OCSP
responses to hold the OCSP-signing usage, and the reference implementation
skips a CA-signed response (Table~\ref{tab:validation}), so the column uses a
delegated responder, which carries its own certificate:
\NocspDelDSAsixtyfive{} bytes under ML-DSA-65 against
\NocspCADSAsixtyfive{} for the CA-signed form.

The timestamp multiplies the certificate term rather than adding to it. Under
ES256 it takes the manifest from \NtsSigES{} to \NtsTstES{} bytes; under
ML-DSA-65, from \NtsSigDSAsixtyfive{} to \NtsTstDSAsixtyfive{}, a factor of
\NtsGrowthDSAsixtyfive{}, because the time-stamping authority's certificates
cost what the signer's do. Figure~\ref{fig:compose} breaks a timestamped
manifest down by role: under every post-quantum parameter set, certificates
are the largest part and the two manifest signatures a minority.

\begin{figure}[t]
\centering
\includegraphics[width=\linewidth]{figures/compose}
\Description{Horizontal stacked bars, one per algorithm from Ed25519 to
ML-DSA-87, showing the bytes of a timestamped manifest split into manifest
signature, signer leaf certificate, signer intermediate certificate, TSA
signature, TSA certificates, and structure. The bars grow from \NtsTstkBEd{} kB for
Ed25519 to \NtsTstkBDSAeightyseven{} kB for ML-DSA-87, and certificates are the largest
segments in the post-quantum bars.}
\caption{What a timestamped manifest is made of, by role: the \emph{+ timestamp} column of Table~\ref{tab:timestamp}, roots omitted, without the OCSP response. Structure covers COSE and JUMBF
wrapping, the claim, and the CMS and TSTInfo encoding of the timestamp.}
\label{fig:compose}
\end{figure}

\subsection{What signing actually costs}
\label{sec:hash}

A content credential is not only a signature. The claim contains a hard
binding, a hash over the content, computed before anything is signed. On a
4\,MB phone photograph SHA-256 takes \NhashPhone\,ms on the M1 Pro, so the hash
is \NhashSharePhone\% of hashing and ML-DSA-65 signing together; on a Kodak
image it is \NhashShareKodak\% (Table~\ref{tab:hash} in the supplement). The
platform has hardware SHA-256 and no hardware support for lattice arithmetic,
the configuration least favourable to this conclusion. On images large enough
to be worth signing, the post-quantum signature costs less than the hash the
format already requires; the signer's cost of migrating is bytes, and time
only at small sizes and high rates, which Section~\ref{sec:edge} measures.

\subsection{Beyond ML-DSA}
\label{sec:algsweep}

\begin{table}[t]
\caption{Signature schemes at hierarchy depth three, leaf and intermediate carried. \emph{Chain} is the share of the signed structure taken by certificates; \emph{Embed.} and \emph{Ref.} are the manifest with the chain and with a thumbprint, as a percentage of the median image. Every certificate is real; Falcon carries a placeholder object identifier and its padded signature length. A dash means no COSE identifier is assigned. Signing times through liboqs on Apple M1 Pro, over the claim alone (Section~\ref{sec:timing}). Times are medians over 10 independent runs; their intervals across runs are in the released data.}
\label{tab:algsweep}
\small
\setlength{\tabcolsep}{4pt}
\begin{tabular}{llrrrrrr}
\toprule
 & COSE & Sig. & Chain & Chain & Embed. & Ref. & Sign\\
Algorithm & id & (B) & (B) & (\%) & (\% img) & (\% img) & (ms)\\
\midrule
Ed25519 & -8 & 64 & 1,387 & 94.2 & 2.4 & 1.1 & --\\
ES256 & -7 & 64 & 1,499 & 94.6 & 2.5 & 1.1 & --\\
PS256 & -37 & 256 & 2,291 & 89.1 & 3.4 & 1.3 & --\\
\midrule
ML-DSA-44 & -48 & 2,420 & 8,723 & 78.1 & 11.4 & 3.3 & 0.061\\
ML-DSA-65 & -49 & 3,309 & 11,781 & 78.0 & 15.0 & 4.1 & 0.096\\
ML-DSA-87 & -50 & 4,627 & 15,697 & 77.2 & 19.9 & 5.3 & 0.138\\
Falcon-512 & -- & 666 & 4,372 & 86.4 & 5.7 & 1.7 & 0.155\\
Falcon-1024 & -- & 1,280 & 7,398 & 85.0 & 9.1 & 2.2 & 0.304\\
SLH-DSA-SHA2-128f & -- & 17,088 & 35,558 & 67.5 & 49.9 & 16.9 & 21\\
SLH-DSA-SHA2-128s & -- & 7,856 & 17,094 & 68.5 & 24.2 & 8.3 & 429\\
SLH-DSA-SHA2-192s & -- & 16,224 & 33,862 & 67.6 & 47.5 & 16.1 & 757\\
SLH-DSA-SHA2-256s & -- & 29,792 & 61,030 & 67.2 & 85.4 & 28.7 & 672\\
\bottomrule
\end{tabular}
\end{table}


Table~\ref{tab:algsweep} holds the chain share between \NsweepShareLo\% and
\NsweepShareHi\% for signatures whose sizes span a factor of
\NsweepSigFactor. SLH-DSA is unsuitable for per-image provenance at the
burst and frame rates this paper considers: SHA2-256s produces a manifest of
\NslhBigPct\% of the median image, and the compact variants sign in hundreds
of milliseconds; a workflow that signs a few documents an hour could still
use it. Falcon-512 is close to classical cost, a manifest of \NfalconEmbPct\%
of the median image with the chain embedded, but FN-DSA is unpublished, its
arithmetic is being changed from floating to fixed point, and the timings here
are for the submitted algorithm~\cite{nist_fips206_status,nist_fips206_plan}.

Two controls a reader will ask for are in the supplement. Composite
signatures, built from their two components with real certificates
(Table~\ref{tab:composite}), cost the sum of their parts: an
ML-DSA-65 and ES256 composite is \NcompositeCose{} bytes of COSE against
\NcoseDSAsixtyfive{} for ML-DSA-65 alone, so hedging against a lattice
break costs a few per cent of the manifest. General-purpose compression
(Table~\ref{tab:compress}) removes \NcompressDSAsixtyfive\% of an ML-DSA-65
signature structure and \NcompressEd\% of an Ed25519 one: what compresses is
X.509 names and structure, and post-quantum keys and signatures do not.

\subsection{Chain depth}
\label{sec:chaindepth}

Each further level of the hierarchy adds one certificate to \texttt{x5chain}
and costs the same again (Table~\ref{tab:chaindepth} in the supplement).
Under ML-DSA-65 the difference between a two-level and a four-level hierarchy
is \NdepthSpread{} percentage points of a Kodak image, the order of the
difference between ML-DSA-44 and ML-DSA-87. Leaving a root out of
\texttt{x5chain} that the verifier already holds changes no trust decision;
shortening the hierarchy itself is a change of PKI design, and its security
consequences belong to the operator rather than to this measurement.
